% 第一阶段修订：按数学论证分章；本轮不编译。
% 参考 arXiv:2605.00301v1 的单栏、紧凑摘要、居中节标题与左侧公式编号。
\documentclass[UTF8,11pt,a4paper,scheme=plain,leqno]{ctexart}
\usepackage[margin=2.7cm]{geometry}
\usepackage{amsmath,amssymb,amsthm,mathtools,enumitem,booktabs,longtable,array}
\usepackage{tikz}
\usetikzlibrary{arrows.meta}
\usepackage{xurl}
\usepackage[hidelinks]{hyperref}
\urlstyle{tt}
\ctexset{section={format=\centering\large\bfseries,beforeskip=1.8ex plus .4ex,afterskip=1.2ex},
  subsection={format=\normalsize\bfseries,beforeskip=1.4ex,afterskip=.6ex}}
\setlength{\parindent}{2em}
\setlength{\parskip}{0pt}
\setlength{\emergencystretch}{1.5em}
\numberwithin{equation}{section}
\newtheoremstyle{chineseplain}{.7\baselineskip}{.5\baselineskip}{\normalfont}{}{\bfseries}{.}{.5em}{}
\theoremstyle{chineseplain}
\newtheorem{theorem}{定理}[section]
\newtheorem{lemma}[theorem]{引理}
\newtheorem{remark}[theorem]{说明}
\newtheorem{proposition}[theorem]{命题}
\newcommand{\N}{\mathbb N}
\newcommand{\limin}{\liminf_{x\to\infty}}
\newcommand{\limsu}{\limsup_{x\to\infty}}
\title{\large\bfseries 正下对数密度与无限整除链\\JSP-001022 的数学证明与形式化对应}
\author{}
\date{}
\begin{document}
\maketitle
\begin{abstract}
本文整理正下对数密度整数序列中无限整除链的证明结构。首先证明一般非负数列的有限 Abel 下界，再通过离散双对数主项估计建立密度转移。随后将已有整除链定理分解为环境链构造与集合内子链提取，并说明端点、枚举和极限值域的转换。最后统一列出固定版本 Lean 源码中的声明对应与显式引用关系。本阶段完整展开密度转换与子链提取；随机链构造和矩估计的详细阐释留待后续补充。
\end{abstract}
\section{引言}
整除偏序中的密度条件如何保证一条无限链，是 JSP-001022（Erdős Problem \#1217）的核心问题。本文先讨论有限 Abel 求和所提供的密度转移，再说明已有整除链定理的两个组成部分，最后统一处理形式化中的表示和依赖关系。整除链定理的数学来源为 Alexeev 等人的论文\cite{alexeev2026}；本文是证明阐释，不主张新的数论结论。

\subsection{问题陈述与记号}
考虑严格递增的正整数序列 $a_1<a_2<\cdots$，令 $A=\{a_i:i\ge1\}$。对于 $x>e$，定义
\begin{align}
H_A(x)&=\sum_{a\in A,\,a<x}\frac1a,\\
W_A(x)&=\sum_{a\in A,\,2\le a<x}\frac1{a\log a}.
\end{align}
采用自然对数及严格截断 $a<x$。在 $a=1$ 处不计算 $1/(a\log a)$，等价于将该权重定义为零。记
\begin{equation}\label{eq:densities}
\underline\delta_{\log}(A)=\limin\frac{H_A(x)}{\log x},\qquad
\overline\delta_{\log\log}(A)=\limsu\frac{W_A(x)}{\log\log x}.
\end{equation}
以下统一记 $\Delta(A):=\overline\delta_{\log\log}(A)$。调和和的积分比较及第~\ref{sec:density}~节的全体整数加权和估计给出
$0\le\underline\delta_{\log}(A)\le1$ 和 $0\le\Delta(A)\le1$，故这两个密度均为有限实数；后文链计数的归一化上极限则允许为 $+\infty$。对于一个严格递增的指标序列 $k_1<k_2<\cdots$，定义 $b_i=a_{k_i}$ 和
\begin{equation}
C_{a,k}(x)=\#\{i\ge1:a_{k_i}<x\}.
\end{equation}
\begin{theorem}[JSP-001022 的枚举形式]\label{thm:main}
若
\begin{equation}\label{eq:assumption}
\underline\delta_{\log}(A)>0,
\end{equation}
则存在严格递增的正整数指标序列 $(k_i)_{i\ge1}$，满足
\begin{equation}\label{eq:divides}
a_{k_i}\mid a_{k_{i+1}}\qquad(i\ge1),
\end{equation}
并且
\begin{equation}\label{eq:target}
\limsu\frac{C_{a,k}(x)}{\log\log x}
\ge\limsu\frac{W_A(x)}{\log\log x}.
\end{equation}
式 \eqref{eq:target} 按照允许 $+\infty$ 的扩展非负实数上极限解释。
\end{theorem}
应注意量词顺序：只选择\emph{一条固定}的链，而不是令其随截断参数 $x$ 改变。上极限的不等式也不是在每个足够大的 $x$ 逐点成立的结论。


\paragraph{历史文献中的密度表述。}
原文\cite[p.~431]{erdos1966}的文字使用“正下对数密度”，但式~(1) 印为 $\limsup$；第~432~页式~(5) 的结论则比较两个上极限。本文遵循 JSP-001022 目录\cite{jspcatalog}及 Erd\H{o}s Problem \#1217\cite{erdos1217}的下密度表述，以式~\eqref{eq:assumption}~的 $\liminf$ 为假设。这里仅记录原文的文字与公式差异，不将正上密度与正下密度视为等价，也不更改本文主定理。

\subsection{证明路线与本文范围}
证明的主干为
\[
 \underline\delta_{\log}(A)>0
 \ \Longrightarrow\ \overline\delta_{\log\log}(A)>0
 \ \Longrightarrow\ \text{具有定量访问下界的环境整除链}
 \ \Longrightarrow\ \text{集合内无限子链}.
\]
第~\ref{sec:density}~节展开第一步的有限不等式、误差估计和极限过程。第~\ref{sec:chain}~节从不变权重和随机链构造出发，证明访问恒等式、矩估计、确定性路径选择及子链提取。第~\ref{sec:conversion}~节完成端点与枚举转换。最后一节集中给出固定版本源码的对应关系与核查边界。

\section{有限 Abel 估计与密度转换}\label{sec:density}
这一节把正下对数密度转化为双对数加权和的增长下界。为与有限求和一致，暂记
\[
 H_A^{\le}(x)=\sum_{\substack{n\in A\\1\le n\le x}}\frac1n,
 \qquad W_A^{\le}(x)=\sum_{\substack{n\in A\\2\le n\le x}}\frac1{n\log n}.
\]
它们与引言中严格截断的和至多相差一个边界项，故相应的归一化极限相同。端点转换的详细说明放在第~\ref{sec:conversion}~节。

\subsection{有限区间上的加权不等式}
\begin{lemma}[下极限的最终下界]\label{lem:eventual}
若 $\delta=\underline\delta_{\log}(A)>0$，则对每个 $0<c<\delta$，存在整数 $M\ge2$，使
\begin{equation}\label{eq:lowerH}
 H_A^{\le}(n)\ge c\log n\qquad(n\ge M).
\end{equation}
\end{lemma}
\begin{proof}
由下极限的定义，$H_A^{\le}(x)/\log x>c$ 对充分大的实数 $x$ 成立。取超过该阈值且不小于 $2$ 的整数 $M$，再将 $x$ 限制为整数即可。
\end{proof}

\begin{lemma}[有限 Abel 下界]\label{lem:abel}
设 $u_n\ge0$，$c>0$，$2\le M\le N$，并令
\[
 S(t)=\sum_{k=0}^{t}u_k,\qquad B_M=S(M-1),\qquad
 D_M(t)=\sum_{q=M+1}^{t}\frac{\log q-\log(q-1)}{\log q}.
\]
若 $S(n)\ge c\log n$ 对所有 $M\le n\le N$ 成立，则
\begin{equation}\label{eq:finite-abel}
 cD_M(N)-\frac{B_M}{\log M}
 \le \sum_{n=2}^{N}\frac{u_n}{\log n}.
\end{equation}
这里的前缀下界只在有限区间 $[M,N]$ 上假设。
\end{lemma}
\begin{proof}
对 $w_n=1/\log n$ 作离散分部求和，得到精确恒等式
\begin{equation}\label{eq:abel}
 \sum_{n=M}^{t}\frac{u_n}{\log n}
 =\frac{S(t)}{\log t}-\frac{B_M}{\log M}
 +\sum_{n=M}^{t-1}S(n)
       \left(\frac1{\log n}-\frac1{\log(n+1)}\right)
 \qquad(M\le t\le N).
\end{equation}
括号内的系数非负，故最后一项至少为
\[
 c\sum_{n=M}^{t-1}\log n
       \left(\frac1{\log n}-\frac1{\log(n+1)}\right)
 =cD_M(t).
\]
因而实际上有更强的不等式
\begin{equation}\label{eq:strong-abel}
 cD_M(t)-\frac{B_M}{\log M}+\frac{S(t)}{\log t}
 \le\sum_{n=M}^{t}\frac{u_n}{\log n}.
\end{equation}
取 $t=N$，舍去非负项 $S(N)/\log N$，再补入右侧 $2\le n<M$ 的非负项，即得结论。
\end{proof}

取 $u_0=0$，并对 $n\ge1$ 取 $u_n=\mathbf1_A(n)/n$。此时 $S(n)=H_A^{\le}(n)$。结合引理~\ref{lem:eventual}，式~\eqref{eq:finite-abel}~给出
\begin{equation}\label{eq:weight-specialization}
 W_A^{\le}(N)\ge cD_M(N)-C_0,\qquad
 C_0=\frac1{\log M}\sum_{\substack{1\le n<M\\n\in A}}\frac1n.
\end{equation}
常数 $C_0$ 保留了有限前缀的损失；在该专门化中，$n=1$ 的倒数项仍被计入前缀。

\subsection{离散双对数主项}
下面说明式~\eqref{eq:weight-specialization}~中的离散和为何具有所需的增长速度。
\begin{lemma}\label{lem:main-term}
存在绝对常数 $K\ge0$，使每个整数 $N\ge2$ 满足
\begin{equation}\label{eq:main-term}
 \left|\sum_{q=2}^{N}\frac{\log q-\log(q-1)}{\log q}
       -\log\log N\right|\le K.
\end{equation}
特别地，固定 $M\ge2$ 时，$D_M(N)=\log\log N+O_M(1)$。
\end{lemma}
\begin{proof}
对 $0<a\le b$，积分比较给出
\[
 0\le\log b-\log a-\frac{b-a}{b}
 =\int_a^b\left(\frac1t-\frac1b\right)\,dt
 \le\frac{(b-a)^2}{ab}.
\]
令 $a=\log r$、$b=\log(r+1)$，其中 $r\ge2$。利用
$\log(r+1)-\log r\le1/r$，可得
\begin{equation}\label{eq:increment-error}
\begin{split}
 0\le{}&\log\log(r+1)-\log\log r
       -\frac{\log(r+1)-\log r}{\log(r+1)}\\
 \le{}&\frac1{r^2(\log r)^2}
 \le\frac1{r(\log r)^2}.
\end{split}
\end{equation}
最后的控制级数收敛。将 $r=2,\ldots,N-1$ 的误差求和，双对数增量望远镜相消；另将 $q=2$ 时等于 $1$ 的首项计入常数，便得到~\eqref{eq:main-term}。去掉至 $M$ 为止的固定前缀，即得关于 $D_M$ 的结论。
\end{proof}

记 $E(M)=\sum_{q=2}^{M}(\log q-\log(q-1))/\log q$。由
$D_M(N)=E(N)-E(M)$，式~\eqref{eq:weight-specialization}~具体给出
\begin{equation}\label{eq:natural-bound}
 W_A^{\le}(N)\ge c\log\log N-C\qquad(N\ge M),
 \qquad C=C_0+c\bigl(K+E(M)\bigr).
\end{equation}
因此误差常数与终端截断 $N$ 无关。这正是由有限不等式通向密度结论所需的统一性。

\subsection{从整数截断到实变量}
\begin{proposition}[密度桥梁]\label{prop:bridge}
若 $\underline\delta_{\log}(A)>0$，则
\begin{equation}\label{eq:bridge}
 \underline\delta_{\log}(A)
 \le\limin\frac{W_A(x)}{\log\log x}
 \le\overline\delta_{\log\log}(A)\le1.
\end{equation}
特别地，$\overline\delta_{\log\log}(A)>0$。
\end{proposition}
\begin{proof}
固定 $0<c<\underline\delta_{\log}(A)$。对充分大的实数 $x$，令 $N=\lfloor x\rfloor$，则
$W_A^{\le}(x)=W_A^{\le}(N)$，且
\[
 \frac{\log\log\lfloor x\rfloor}{\log\log x}\longrightarrow1.
\]
具体地，当 $x\ge3$ 时，$2\le N\le x<2N$，故
\[
 0\le\log\log x-\log\log N
 =\log\left(1+\frac{\log(x/N)}{\log N}\right)
 \le\frac{\log2}{\log N}\longrightarrow0.
\]
结合 $\log\log x\to\infty$ 即得上述比值极限。将式~\eqref{eq:natural-bound}~除以 $\log\log x$，便有
\begin{equation}\label{eq:quantbridge}
 \limin\frac{W_A^{\le}(x)}{\log\log x}\ge c.
\end{equation}
另一方面，对递减函数 $1/(t\log t)$ 作积分比较，得到
\[
 0\le W_A^{\le}(x)\le\sum_{2\le n\le x}\frac1{n\log n}
 =\log\log x+O(1).
\]
这同时说明加权密度有限且至多为 $1$。改变单个边界项不改变归一化极限，故~\eqref{eq:quantbridge}~也适用于 $W_A(x)$。最后令 $c\uparrow\underline\delta_{\log}(A)$ 即得~\eqref{eq:bridge}。
\end{proof}

本节给出的是数学层面的常数比较。扩展源码中密度桥梁的对外结论仅为正性蕴含；证明内部使用固定正数及最终下界 $c/2$，并不声明已经形式化了~\eqref{eq:bridge}~的全部强化形式。已有 plby 形式化则明确输出下对数密度不超过双对数加权上密度的比较\cite{plby}。这些接口的差别及离散主项估计的来源，统一见最后一节。

\section{环境整除链与集合内子链}\label{sec:chain}
密度转换之后，还需要一个实质性的数论结果。Alexeev 等人\cite[Theorem 1.6]{alexeev2026}证明了如下结论；其概率论证位于该文第~9~节。
\begin{theorem}[双对数密度整除链定理，已有结果]\label{thm:upstream}
设 $A\subseteq\N_{>0}$，且
\[
 \Delta(A)=\limsu\frac{W_A^{\le}(x)}{\log\log x}>0.
\]
则存在严格递增的无限整除链 $b_0\mid b_1\mid b_2\mid\cdots$，其所有元素属于 $A$，并满足
\begin{equation}\label{eq:upstream-bound}
 \limsu\frac{\#\{j\ge0:b_j\le x\}}{\log\log x}\ge\Delta(A).
\end{equation}
左侧允许取值 $+\infty$。
\end{theorem}

\subsection{von Mangoldt 函数与不变权重}
以下状态空间为 $S=\N_{>0}$，时间指标为 $\N_0$。记
\[
 \Lambda(n)=\begin{cases}\log p,&n=p^k,\ p\text{ 为素数},\ k\ge1,\\0,&\text{其他情形}.
 \end{cases}
\]
特别地，$\Lambda(1)=0$。若 $n=\prod_p p^{v_p(n)}$，则
\begin{equation}\label{eq:mangoldt-divisors}
 \sum_{d\mid n}\Lambda(d)=\sum_p\sum_{k=1}^{v_p(n)}\log p
 =\sum_p v_p(n)\log p=\log n.
\end{equation}
$n=1$ 时两边均为零。

令 $\zeta(s)=\sum_{r\ge1}r^{-s}$，其中 $s>1$。原论文\cite[Example 2.8]{alexeev2026}和固定版本 Prim 使用的权重是
\begin{equation}\label{eq:mangoldt-weight-exact}
 \boxed{\nu_\Lambda(1)=1,\qquad
 \nu_\Lambda(n)=\int_1^\infty\frac{\log n}{\zeta(s)n^s}\,ds\quad(n\ge2).}
\end{equation}
Prim 的积分区间为 $(1,\infty)$，分母写作复 Riemann $\zeta$ 函数在实参数处的实部；在此区间它就是上式的正实数 $\zeta(s)$。本节始终排除状态 $0$。
因为 $\zeta(s)>1$，被积函数连续且严格为正，并由可积函数 $\log n\,n^{-s}$ 控制，所以
\[
 0<\nu_\Lambda(n)\le\int_1^\infty\log n\,n^{-s}\,ds=\frac1n
 \qquad(n\ge2).
\]
这同时说明积分的有限性与权重的严格正性。Erd\H{o}s 权重则为
$\nu_0(n)=1/(n\log n)$（$n\ge2$）；为统一求和可令 $\nu_0(1)=0$。两种权重在 $1$ 处尤其不能混淆，其总和的比较留待第~3.3~小节。

不变权重的作用是平衡从较大整数流入较小整数的\emph{严格整除}转移。它不是总质量为 $1$ 的概率分布，也不是把状态 $1$ 的吸收自环包括在内的平稳分布。所需的准确恒等式如下。
\begin{proposition}[严格入流恒等式]\label{prop:mangoldt-inflow}
对每个 $n\ge1$，均有收敛的非负级数恒等式
\begin{equation}\label{eq:mangoldt-inflow}
 \boxed{\sum_{q\ge2}\nu_\Lambda(nq)\frac{\Lambda(q)}{\log(nq)}
 =\nu_\Lambda(n).}
\end{equation}
\end{proposition}
\begin{proof}
写 $h(s)=1/\zeta(s)$。由积分比较，
\[
 \frac1{s-1}\le\zeta(s)\le1+\frac1{s-1},\qquad \zeta(s)\ge1,
\]
故 $h(1+)=0$，$h(+\infty)=1$。Euler 乘积的对数导数给出
\begin{equation}\label{eq:reciprocal-zeta-derivative}
 -\frac{\zeta'(s)}{\zeta(s)}=\sum_{q\ge2}\frac{\Lambda(q)}{q^s},
 \qquad h'(s)=h(s)\sum_{q\ge2}\frac{\Lambda(q)}{q^s}\ge0.
\end{equation}
这里的逐项求导可在每个 $s\ge1+\varepsilon$ 的区间上用绝对、一致收敛证明；例如 $\Lambda(q)\le\log q$，而 $\sum_{q\ge2}(\log q)q^{-1-\varepsilon}$ 收敛。
在有限区间应用微积分基本定理，再令端点趋向 $1$ 与无穷，由非负性得到
\begin{equation}\label{eq:reciprocal-zeta-total-derivative}
 \int_1^\infty h'(s)\,ds=1.
\end{equation}
将~\eqref{eq:mangoldt-weight-exact}~代入左端，由 Tonelli 定理交换非负求和与积分，得到
\begin{align*}
 \sum_{q\ge2}\nu_\Lambda(nq)\frac{\Lambda(q)}{\log(nq)}
 &=\int_1^\infty n^{-s}h(s)\sum_{q\ge2}\Lambda(q)q^{-s}\,ds\\
 &=\int_1^\infty n^{-s}h'(s)\,ds.
\end{align*}
交换时可先允许积分取扩展非负实数值；最后一个积分不超过 $\int h'=1$，因此也证明了级数的有限性，没有预先假设待证的可和性。
当 $n=1$ 时，结果由~\eqref{eq:reciprocal-zeta-total-derivative}~给出，恰为定义中的 $\nu_\Lambda(1)=1$。
当 $n\ge2$ 时，在 $[a,b]\subset(1,\infty)$ 上分部积分，然后令 $a\downarrow1$、$b\to\infty$：
\[
 \int_1^\infty n^{-s}h'(s)\,ds
 =[n^{-s}h(s)]_{1+}^{\infty}
   +\log n\int_1^\infty n^{-s}h(s)\,ds
 =\nu_\Lambda(n).
\]
边界项均为零；左侧由 $h'$ 控制，右侧被积函数由 $n^{-s}$ 控制，所以此极限步骤合法。
\end{proof}

上述推导对应 Prim 的不变递推声明。正文采用非负积分的 Tonelli 论证；源码则先建立绝对可积求和的桥接等式，再调用倒数 $\zeta$ 的端点评估。两者的数学职责相同，但不能据此声称正文的每个中间不等式都已有同名 Lean 声明。

\subsection{从下降整除链构造向上转移核}
\paragraph{下降核与边界。}
在 $S$ 上定义
\begin{equation}\label{eq:downward-kernel}
 P_\downarrow(n,m)=
 \begin{cases}
 1,&n=m=1,\\
 \Lambda(n/m)/\log n,&n\ge2,\ m\mid n,\ m<n,\\
 0,&\text{其他情形}.
 \end{cases}
\end{equation}
$n\ge2$ 时分母严格为正，故所有转移概率非负。映射 $m\mapsto q=n/m$ 将正真因子与满足 $q\mid n$、$q>1$ 的因子一一对应。由~\eqref{eq:mangoldt-divisors}，
\[
 \sum_{m\in S}P_\downarrow(n,m)
 =\frac1{\log n}\sum_{\substack{q\mid n\\q>1}}\Lambda(q)=1.
\]
$n=1$ 的行和则由吸收自环给出。因而它是一个下降 Markov 核，除吸收状态 $1$ 外仅沿严格整除关系转移。

\paragraph{伴随向上核。}
利用权重的严格正性，定义
\begin{equation}\label{eq:upward-kernel}
 U(m,n)=P_\uparrow(m,n)=
 \begin{cases}
 \displaystyle\frac{\nu_\Lambda(n)}{\nu_\Lambda(m)}P_\downarrow(n,m),&m\ne n,\\
 0,&m=n.
 \end{cases}
\end{equation}
等价地，$U(m,mq)=\nu_\Lambda(mq)\Lambda(q)/(\nu_\Lambda(m)\log(mq))$（$q\ge2$），其余项为零。因此 $U\ge0$，且
\[
 U(m,n)>0\ \Longrightarrow\ m<n,\quad m\mid n;
\]
更精确地，商 $n/m$ 必须为素数幂。重新以 $n=mq$ 编号，再应用命题~\ref{prop:mangoldt-inflow}，得到每个 $m\ge1$ 的行和
\begin{equation}\label{eq:upward-normalization}
 \sum_{n\in S}U(m,n)
 =\frac1{\nu_\Lambda(m)}\sum_{q\ge2}
    \nu_\Lambda(mq)\frac{\Lambda(q)}{\log(mq)}=1.
\end{equation}
伴随平衡关系
\[
 \nu_\Lambda(m)U(m,n)=\nu_\Lambda(n)P_\downarrow(n,m)
\]
仅对 $m\ne n$ 声明。在 $(1,1)$ 处，右侧为 $1$ 而左侧为 $0$；若把下降核的吸收自环也搬到向上核，则会破坏归一化。

原论文\cite[Definition 2.9, Example 2.11]{alexeev2026}允许加入状态 $\infty$，并把有限行和的亏损赋给 $U(m,\infty)$。本例由~\eqref{eq:upward-normalization}~知该亏损为零。从 $1$ 出发，在任一有限步到达 $\infty$ 的概率为零，所有步的可数并仍为零。因此可直接在 $S$ 上构造无限时间的链；整数状态随时间趋于无穷，并不意味着某一有限时刻跳到附加状态 $\infty$。

\paragraph{同一无限路径空间上的构造。}
给 $S$ 配备离散 $\sigma$-代数，以 $\delta_1$ 为初始分布，以 $U$ 为转移核。每一行都是概率质量函数，可测性在可数离散空间上自动成立。Ionescu--Tulcea 定理于是给出乘积可测空间 $S^{\N_0}$ 上的唯一概率测度，使坐标过程 $(X_i)$ 的柱集概率为
\begin{equation}\label{eq:chain-cylinders}
 \mathbb P(X_0=n_0,\ldots,X_r=n_r)
 =\mathbf1_{\{n_0=1\}}\prod_{i=0}^{r-1}U(n_i,n_{i+1}).
\end{equation}
这些有限维分布的相容性正是行和为 $1$：对最后一个坐标求和就恢复前一阶分布。该定理把全部相容分布放到\emph{同一个}无限乘积空间上，坐标投影均可测。特别地，对任意 $k,m,n$ 有
\begin{equation}\label{eq:chain-pair-law}
 \mathbb P(X_k=m,X_{k+1}=n)=\mathbb P(X_k=m)U(m,n).
\end{equation}

还需将几乎处处成立的支撑条件整理成逐路径成立的条件。令
\[
 G=\{X_0=1\}\cap\bigcap_{k\ge0}\{U(X_k,X_{k+1})>0\}.
\]
对固定 $k$，零转移事件是满足 $U(m,n)=0$ 的可数多个二坐标柱集的并；每个柱集由~\eqref{eq:chain-pair-law}~知概率为零。因此 $G$ 可测且概率为 $1$。把概率空间限制到 $G$ 不改变任何概率，也不需重新归一化。在这个空间上，\emph{每条}路径都满足
\[
 X_0=1,\qquad X_k<X_{k+1},\qquad X_k\mid X_{k+1}\quad(k\ge0).
\]
特别地，$X_{k+1}\ge2X_k$，从而 $X_k\ge2^k$，得到无限的严格递增整除链。这一辅助增长估计也保证任何有限截断内的访问次数有限。

Prim 将核写在全部自然数上，原始 $U$ 的第 $0$ 行为零，并不归一化。构造概率核时，它人为把该行补为 $\delta_1$；从初始状态 $1$ 出发的支撑归纳保证永远不访问 $0$，故补行不影响上述正整数过程。随后源码使用路径测度及一个满测度子类型，使严格递增和整除性质对该子类型的每个样本成立。
其导出的 Markov 接口是~\eqref{eq:chain-pair-law}~这一相邻二坐标恒等式；完整的无限路径分布来自实际的路径测度构造，不能仅凭该接口名称推断全部条件概率性质。

至此完成的是无限随机环境链的构造。要从中选出对给定集合 $A$ 达到密度下界的确定性样本，仍须后续的访问恒等式、矩估计和极限论证。

\subsection{访问概率与加权和：待展开}
下一步需要证明
\begin{equation}\label{eq:visit-goal}
 \sum_{i\ge0}\mathbb P(X_i=n)=\nu_\Lambda(n),
 \qquad
 \sum_{\substack{n\in A\\n\le x}}\nu_\Lambda(n)=W_A^{\le}(x)+O(1).
\end{equation}
这里权重比较的误差为何一致有界尚待展开，不能仅由本节已证的 $\nu_\Lambda(n)\le1/n$ 推出。

\subsection{一阶矩和二阶矩：待展开}
记 $Z_A(x)=\#\{i\ge0:X_i\in A,\ X_i\le x\}$。后续需要论证
\begin{equation}\label{eq:moment-goals}
 \mathbb E Z_A(x)=\sum_{\substack{n\in A\\n\le x}}\nu_\Lambda(n),
 \qquad \mathbb E[Z_A(x)^2]\ll(\log\log x)^2.
\end{equation}
本轮不展开质因子个数函数与二阶矩的关系。

\subsection{反向 Fatou 与确定性路径：待展开}
沿使加权密度趋于 $\Delta(A)$ 的尺度序列归一化后，还需利用矩控制与反向 Fatou 论证，选出\emph{同一条}确定性环境链 $m_0\mid m_1\mid\cdots$，满足
\begin{equation}\label{eq:ambient-hits}
 \limsu\frac{\#\{i\ge0:m_i\in A,\ m_i\le x\}}{\log\log x}
 \ge\Delta(A).
\end{equation}
本节构造的无限路径空间是这一步的前提；仅在每个尺度分别选出一条有限链，仍不能保证这些链相容。

\subsection{从访问指标抽取子链}
\begin{lemma}[集合内子链提取]\label{lem:extract}
若 $\Delta(A)>0$，且严格递增的环境整除链 $(m_i)$ 满足~\eqref{eq:ambient-hits}，则它包含满足~\eqref{eq:upstream-bound}~的集合内无限子链。
\end{lemma}
\begin{proof}
设 $I=\{i\ge0:m_i\in A\}$。若 $I$ 有限，则访问次数至多为 $\#I$，除以趋于无穷的 $\log\log x$ 后极限为零，与~\eqref{eq:ambient-hits}~矛盾。故可将 $I$ 递增枚举为 $i_0<i_1<\cdots$，并令 $b_j=m_{i_j}$。

严格递增性立即保留。由环境链相邻项的整除关系及传递性，$m_r\mid m_s$ 对所有 $r<s$ 成立，故 $b_j\mid b_{j+1}$。又因 $j\mapsto i_j$ 是到 $I$ 的双射，对每个 $x$ 都有精确的计数恒等式
\[
 \#\{j\ge0:b_j\le x\}
 =\#\{i\ge0:m_i\in A,\ m_i\le x\}.
\]
因而所需的上极限下界原样保留。
\end{proof}

这样，已有整除链定理的证明职责被分解为：先证明概率方法给出的~\eqref{eq:ambient-hits}，再应用引理~\ref{lem:extract}。本轮已展开随机环境链本身的构造；从随机链得到上述密度下界的步骤仍作有明确来源的引用，子链提取则保留完整推导。尚不把这一阶段描述为对上游全部证明的逐项复现。

\section{截断、枚举与主定理的组合}\label{sec:conversion}
本节说明如何把集合内整除链结论转成引言中的正整数枚举形式。这里先处理数学对象，形式化声明的对应关系留到最后一节。

\subsection{端点与极限值域}
严格截断 $n<x$ 和非严格截断 $n\le x$ 只在 $x$ 是整数且该整数属于所计集合时有差别。对充分大的 $x$，有
\[
 0\le H_A^{\le}(x)-H_A(x)\le1,\qquad
 0\le W_A^{\le}(x)-W_A(x)\le\frac1{2\log2}.
\]
任意严格递增链的两种截断计数也至多相差 $1$。分别除以 $\log x$ 或 $\log\log x$ 后，这些误差均趋于零。因此倒数和的下极限、加权和的上极限及链计数的上极限都不受端点改变的影响。

加权密度由~\eqref{eq:bridge}~中的全体整数上界控制，是有限实数；链计数的归一化上极限则可能为 $+\infty$。例如 $1,2,4,8,\ldots$ 的计数为 $\log x/\log2+O(1)$。故主定理的计数上极限应放在 $[0,+\infty]$ 中解释。对最终非负且有界的加权比值，嵌入扩展非负实数与取上极限相容；计数一侧不需要额外假设有界。

\subsection{正整数子序列}
\begin{proof}[定理~\ref{thm:main}~的证明（使用已有整除链定理）]
令 $A=\{a_i:i\ge1\}$。严格递增性保证枚举没有重复，所以有限截断下的集合求和等于指标求和。由命题~\ref{prop:bridge}，$\Delta(A)>0$。应用定理~\ref{thm:upstream}，得到 $A$ 中严格递增的链 $(b_j)_{j\ge0}$ 及~\eqref{eq:upstream-bound}。

对每个 $i\ge1$，由 $b_{i-1}\in A$ 及枚举的单射性，存在唯一正整数 $k_i$ 使 $a_{k_i}=b_{i-1}$。如果 $k_i\ge k_{i+1}$，则 $a_{k_i}\ge a_{k_{i+1}}$，与 $b_{i-1}<b_i$ 矛盾。因此 $k_i<k_{i+1}$，且
\[
 a_{k_i}=b_{i-1}\mid b_i=a_{k_{i+1}}.
\]
指标平移 $i\mapsto i-1$ 给出相同的截断计数。再将~\eqref{eq:upstream-bound}~的两侧改为严格截断，即得~\eqref{eq:target}。
\end{proof}

在这条推导中，密度桥梁只负责提供调用上游定理所需的正性。最终计数下界是原集合的完整加权上密度 $\Delta(A)$，而不是桥梁证明中选取的辅助常数 $c$ 或 $c/2$。

\section{Lean 对应关系、源码依赖与核查范围}\label{sec:formal}
\subsection{版本与数学接口}
扩展源码取自提交 \nolinkurl{3992abb235729e5ebd0359d897470d6101acf908}；Prim 取自 \nolinkurl{a303f6bd23bb8f29a833d1d70327528359180995}。下列行号均对应这些固定文件，而非随时可能变化的默认分支。

本节区分数学推导、源码显式引用和编译后证明项依赖。正文的有限 Abel 恒等式给出式~\eqref{eq:strong-abel}；扩展以该强化不等式为归纳不变量，使用自然数归纳实现有限估计。离散主项引理~\ref{lem:main-term}~对应 Prim 已有的 \nolinkurl{mangoldt_log_reciprocal_main_term_bound}（第 11114 行），并非扩展中新证明的上游估计。

扩展的 \nolinkurl{jsp_001022_density_bridge} 只输出密度的正性蕴含。正文给出的更强常数比较不可等同于该声明的类型。plby 的 \nolinkurl{lowerLogDensity_le_weightedRate}（Density 第 533 行）已经形式化了下对数密度与加权上密度的比较；其第 259 行的有限 Abel 引理使用积分形式，与本扩展的离散归纳实现不同。这一实现差别不构成数学优先权或首次形式化的证明。

\subsection{十五个声明及其显式引用}
表中省略公共命名空间 \nolinkurl{LeanCompletion}；除 E08 外，还省略前缀 \nolinkurl{jsp_001022_}。末列只列扩展内部的直接显式引用，空白处表示未发现其他十四个声明的显式引用，不表示没有 Mathlib 或 Prim 依赖。

{\footnotesize
\setlength{\tabcolsep}{3pt}
\begin{longtable}{@{}p{.06\linewidth}p{.37\linewidth}p{.07\linewidth}p{.27\linewidth}p{.15\linewidth}@{}}
\toprule
编号 & 声明名 & 行号 & 数学职责 & 引用\\\midrule
\endfirsthead
\toprule
编号 & 声明名 & 行号 & 数学职责 & 引用\\\midrule
\endhead
E01 & \nolinkurl{eventual_lower_bound} & 7 & 正下极限的最终下界 & ---\\
E02 & \nolinkurl{finite_abel_lower_bound} & 38 & 一般非负数列的有限 Abel 不等式 & ---\\
E03 & \nolinkurl{reciprocal_weight_specialization} & 112 & 倒数权专门化 & E02\\
E04 & \nolinkurl{natural_loglog_lower_bound} & 155 & 自然数截断的双对数下界 & E03\\
E05 & \nolinkurl{eventual_to_finite_prefix} & 198 & 实变量最终界转为有限前缀界 & ---\\
E06 & \nolinkurl{density_to_natural_erdos_lower_bound} & 246 & 合成自然数尺度的加权下界 & E01, E04, E05\\
E07 & \nolinkurl{density_bridge} & 293 & 正下对数密度推出正加权上密度 & E06\\
E08 & \nolinkurl{erdos_1217_explicit} & 428 & 展开 Prim 的整除链定理 & ---\\
E09 & \nolinkurl{chain_strict_cutoff} & 447 & 链计数的严格截断 & ---\\
E10 & \nolinkurl{reciprocal_strict_cutoff} & 528 & 倒数和的严格截断 & ---\\
E11 & \nolinkurl{erdos_strict_cutoff} & 649 & 加权和的严格截断 & ---\\
E12 & \nolinkurl{erdos_density_ennreal} & 868 & 实数与扩展非负实数的加权上极限 & E11\\
E13 & \nolinkurl{positive_enumeration_sums} & 1029 & 一般权重下的枚举与集合求和 & ---\\
E14 & \nolinkurl{positive_subsequence_indices} & 1077 & 正整数子序列及精确计数 & ---\\
E15 & \nolinkurl{original_enumeration_chain} & 1136 & 最终枚举定理 & E07, E08, E09, E10, E12, E13, E14\\
\bottomrule
\end{longtable}
}

E04 直接引用 Prim 的离散主项估计；E07 还引用第 11545 行的取整误差界。E11、E12 同样利用第 11114 行的主项估计取得最终上界。E08 直接调用第 14584 行的 \nolinkurl{erdos_sarkozy_szemeredi_1217}。这解释了为什么密度、端点与数系转换并不是一条没有分支的线性依赖链。

此表由固定源码在去掉注释和字符串后提取声明标识符，并核读主要调用处得到。仓库中的 \nolinkurl{docs/JSP001022_SOURCE_GUIDE.md} 整理了固定源码链接、十五个声明的引用关系及上游阅读顺序。临时索引与下载源码留在本地，不作为仓库交付材料。这里的引用关系属于源码级核读，不包括策略展开、类型类搜索及简化器自动引入的全部依赖，也不等同于编译后证明项的依赖审计。

\subsection{上游概率证明的阅读入口}
Prim 中环境链定理的直接组合关系为：第 14423 行的 \nolinkurl{probabilistic_dense_ambient_chain} 使用权重密度比较（第 7962 行）与路径选取（第 14397 行）；后者使用随机模型存在性（第 14328 行）和反向 Fatou 路径提取（第 14366 行）。最终定理将环境链与第 14464 行的 \nolinkurl{dense_hits_subchain_in_set} 组合。

随机模型存在性的证明体进一步调用构造数据存在性、二阶矩界及反向 Fatou 提取原则。它没有仅凭“随机模型”这个名称就完成概率论证。下一阶段按以下顺序展开其数学内容：
\begin{enumerate}[leftmargin=2em]
\item 不变权重与转移核：第 9423、9677 行；路径测度构造：第 9900 行。
\item 访问恒等式：第 10232 行；权重的可求和误差与密度比较：第 7727、7962 行。
\item 双点访问估计：第 10913 行；质因子数加权和：第 12817 行；二阶矩合成：第 12870、14149 行。
\item 一阶、二阶矩接口：第 13677 行；截断反向 Fatou：第 13831 行；固定路径提取：第 14366 行。
\end{enumerate}
这些入口均已保存到本地源码检索材料；正文目前尚未逐项展开其证明。该待展开范围不同于声称已有 Lean 定理缺少证明。

\subsection{现有审计材料与工作边界}
扩展仓库现有 \nolinkurl{LeanCompletion/Audit.lean} 会独立重述最终类型，并用 \nolinkurl{Expr.getUsedConstants} 导出十五个声明的类型和证明体引用，另有公理检查及模块路径检查。现有 \nolinkurl{verification/targets.json} 保留核查目标，\nolinkurl{docs/VERIFICATION.md} 保留复现步骤。本轮只读取这些文件，没有执行其中的构建或审计命令。

仓库的 2026 年 9 月 28 日验证摘要报告所选证明提交曾通过检查，但也明确说明使用了既有本地依赖，当前树不再分发原始执行日志。因此该历史记录不能写成本轮新的独立复现。本轮只修改 LaTeX、撰写说明并阅读源码，没有安装 Lean、下载依赖环境或运行编译。

主定理沿用当前问题表述的正下对数密度。仓库的陈述对应文档报告 1966 年原文第 431 页文字与所印极限符号存在差异；本轮没有重新核对该历史页面，也不据此宣称正下密度与正上密度等价。对任意无限正整数集合构造递增枚举虽是通常数学事实，当前扩展的最终接口仍是正整数枚举形式，本稿不额外声称已提供集合枚举构造的 Lean 包装。

\begin{thebibliography}{9}
\small
\setlength{\itemsep}{.45\baselineskip}
\bibitem{erdos1966} P. Erd\H{o}s, A. S\'ark\"ozy and E. Szemer\'edi,
\emph{On divisibility properties of sequences of integers}, Studia Sci. Math. Hungar. \textbf{1} (1966), 431--435.
\href{https://users.renyi.hu/~p_erdos/1966-09.pdf}{原文扫描件}，特别是第 431--432 页。
\bibitem{alexeev2026} B. Alexeev, K. Barreto, Y. Li, J. D. Lichtman, L. Price, J. I. Shah, Q. Tang and T. Tao,
\emph{Primitive sets and von Mangoldt chains: Erd\H{o}s Problem \#1196 and beyond},
\href{https://arxiv.org/abs/2605.00301v1}{arXiv:2605.00301v1} (2026)，定理 1.6、第 2、9 节。
\bibitem{prim} YuanheZ, \emph{Prim},
\href{https://github.com/YuanheZ/Prim/tree/a303f6bd23bb8f29a833d1d70327528359180995}{固定源码版本}：
\nolinkurl{a303f6bd23bb8f29a833d1d70327528359180995}。
\bibitem{extension} huam54925-cpu, \emph{jsp-001022-lean-extension},
\href{https://github.com/huam54925-cpu/jsp-001022-lean-extension/tree/3992abb235729e5ebd0359d897470d6101acf908}{选定证明版本}：
\nolinkurl{3992abb235729e5ebd0359d897470d6101acf908}。
\bibitem{plby} plby, \emph{lean-proofs},
\href{https://github.com/plby/lean-proofs/tree/8822f7ddef30fadbd92e1c6ab4ed897af356af5e/src/latest/ErdosProblems/Erdos1217}{Erdos1217 模块}，特别是 Density、Resolution 与 Reindex；固定提交：
\nolinkurl{8822f7ddef30fadbd92e1c6ab4ed897af356af5e}。
\bibitem{jspcatalog} The Justin Sun Prize,
\href{https://github.com/TheJustinSunPrize/awards/blob/8332eaab70094567cb3a2137e33b6ea5a2e316d8/problems/catalog-1001-1022.md\#JSP-001022}{JSP-001022 目录条目}，固定目录版本 \nolinkurl{8332eaab70094567cb3a2137e33b6ea5a2e316d8}；2026 年 10 月 9 日核对的在线目录仍采用正下对数密度。
\bibitem{erdos1217} T. F. Bloom,
\href{https://www.erdosproblems.com/1217}{Erd\H{o}s Problem \#1217}，题目表述采用正下对数密度；2026 年 10 月 9 日经检索索引核对。
\end{thebibliography}

\end{document}
